% D06-001. PDR-05, разделы 4 и 6.
\section{Подготовка массива для проверки измерительных органов}
\label{sec:ch5_population_audit}
\conceptsource{Статья 5, разделы 4 и 6. Перенесены состав
измерительных массивов, единицы усреднения, сравнение органов
без физического эталона и отбор сложных режимов.}

Массовая проверка нужна для отбора осциллограмм, на которых
различия органов имеют инженерный смысл. Если ответы расходятся
при малом токе, следует проверить чувствительность. Если расхождение
возникает при коммутации, нужно рассмотреть смену направления
и устойчивость выхода. Поэтому сначала оценивается состав архива,
затем находятся характерные расхождения и выбираются записи
для просмотра токов, напряжений и дорожек решений.

\subsection{Источники и физические условия}
\conceptsource{Статья 5, таблица 1 и описание двух доменов.
Числа относятся к исследовательскому снимку статьи, а не
к неизменному составу постоянно пополняемого архива.}

Для массового сопоставления использованы два массива
реальных осциллограмм. Распределительный архив Open\_EE
относится к сетям 6--35~кВ. Он ближе к условиям применения
БАВР на промышленных секциях с двигательной нагрузкой.
В обучении доля French/RTE была ограничена. Этот архив полезен
для проверки иных режимов сети, но по составу измерений и процессов
не полностью соответствует основной задаче РНМ для БАВР.
Полный пригодный архив и доля записей в обучении различаются.

Массив French/RTE содержит процессы магистральных сетей
225 и 400~кВ. Различаются параметры сети, нагрузка,
ёмкостные токи линий и условия регистрации.

\begin{table}[htbp]
\centering\small
\caption{Массивы массового сопоставления РНМ}
\label{tab:ch5_population}
\begin{tabular}{lrrr}
\toprule
Источник & Всего записей & Пригодно & Длительность, ч \\
\midrule
Open\_EE & 44\,773 & 36\,737 & 39,571 \\
French/RTE & 12\,053 & 12\,053 & 10,919 \\
Всего & 56\,826 & 48\,790 & 50,490 \\
\bottomrule
\end{tabular}
\end{table}

Пригодность в таблице~\ref{tab:ch5_population} означает
наличие минимального состава двух токов и двух напряжений
с конечными числовыми отсчётами. Объявленный, но фактически
не заполненный канал измерением не считается.
Нулевой числовой сигнал отличается от отсутствующего.
В Open\_EE по этому условию исключено 8036 записей.

В пригодной части третья фаза тока восстанавливалась
в 6763 записях, преимущественно при отсутствии фазы B.
В этой подготовке принималась нулевая сумма фазных токов.
Восстановление третьего фазного напряжения потребовалось
для пяти записей. Ограничения нулевой суммы фазных
величин рассмотрены в главе~\ref{ch:oscillogram_analysis}.
Здесь они существенны и для сравнения органов.
При однофазной наладочной подаче испытательных токов такая подстановка
может создать искусственную противофазную составляющую
и изменить ответ органа прямой последовательности.
Подобные записи рассматриваются отдельно от сетевых повреждений.

Внутреннее представление оставляет запас по току
до двадцати номиналов и по напряжению до трёх номиналов.
Эти масштабы отличаются от относительных единиц,
в которых задаются уставки органов.
Например, ток $0{,}05I_{\text{ном}}$ после дополнительного
деления на 20 уже не имеет численного значения 0,05.
Поэтому при сопоставлении сохраняется полная цепочка
перехода от физических измерений к рабочему критерию.
Масштабирование не заменяет ограничения входа и не требует
обрезать каждый ток выше номинального значения.

Магистральный архив используется для проверки поведения
в другом физическом окружении. При значительной реактивной
составляющей знак активной мощности и угловой критерий
с ненулевым рабочим углом могут различаться.
Такое расхождение требует разбора принятой характеристики.
Оно не означает автоматически ошибку вычислений.
Поскольку часть French/RTE участвует в обучении,
весь этот источник не называется независимым внешним тестом.

\iffalse
(Примечание автора: Стоит отдельно ещё раз сказать, что эта база осциллограмм использовалась в малой доле, так как не совсем подходит...)
(Примечание автора: Ок, закрываем.)
\fi
\authorreviewnote{AA157}


\subsection{Сохраняемые решения и единицы усреднения}
\conceptsource{Статья 5, подразделы 6.1--6.3. Направление,
применимость и запас до границы; статистическая единица
осциллограмма или группа точных копий.}

Для каждого момента сохраняются направление,
применимость, запас до границы срабатывания
и внутренняя диагностическая оценка органа.
Запас позволяет отличать устойчивое решение
от ответа вблизи уставки. Он зависит от формулы
и масштаба величин. Диагностическая оценка также
не становится вероятностью правильного направления
без самостоятельной проверки её соответствия ошибкам.

Статистика рассчитывается несколькими способами.
При пофайловом усреднении сначала определяется доля
состояния внутри каждой записи, затем записи получают
одинаковый вес. При усреднении по длительности
равный вес имеет секунда наблюдения. При подсчёте
по отсчётам больший вклад получают также регистраторы
с более высокой частотой дискретизации.
Различие этих результатов показывает, с чем связана
наблюдаемая особенность массива.

Оценка разброса выполняется повторным выбором целых записей
или групп их копий с возвращением. Для каждого такого набора
заново рассчитывается показатель. Полученный диапазон показывает,
насколько результат чувствителен к составу проверяемого материала.
В статистике этот приём называют бутстрепом.

Соседние решения одного скользящего окна связаны.
Поэтому единицей оценки неопределённости служит
осциллограмма или группа её точных копий.
Повторная выборка выполняется целыми группами,
с сохранением внутренней временной структуры.
Точное совпадение файлов не исчерпывает все зависимости.
Разные регистрации одного повреждения также могут
описывать одно физическое событие.
\ifdefined\AuthorReview
\par\smallskip\textit{Служебное уточнение. В статье для массового
бутстрепа написано 2000 повторов, но прежняя проверка
генератора установила 1000. Здесь спорное число не перенесено.
Конкретные интервалы требуют ссылки на сохранённый расчёт.}\par
\fi

\iffalse
(Примечание автора: Вообще, круто что метрики описываются. Но… Но не понятно, а чё и зачем это вообще? Надо как-то чуть раньше сказать понятным, человеческим языков... Повторюсь, это пишется для электроэнергетиков. Ну и слово бутрстрап они не знают, хотя основное замечание не к слову. Может чуть порядок разделов поменять. В общем с начала цель не очень ясна...)
(Примечание автора: Ок, закрываем.)
\fi
\authorreviewnote{AA158}


\subsection{Расхождения как средство поиска режимов}
\conceptsource{Статья 5, подразделы 6.4--6.7 и постановка
массового анализа. Согласие не подменяет экспертную правильность.}

Первое сопоставление выполняется без предположения,
что какой-либо орган всегда прав.
Для каждой пары выделяются моменты совместной
применимости, совпадения направления и расхождения.
Дополнительно рассматриваются совпадение полного
состояния и доля времени, на которой орган вообще
выдаёт применимый ответ.
Так исключается скрытое повышение согласия за счёт
отказа от решения на сложных участках.

Помимо доли совпадений используются коэффициенты
согласия с поправкой на частоты состояний.
Однако даже полное совпадение нескольких органов
может отражать общую недостаточность их характеристик.
Большинство голосов поэтому служит ориентиром для отбора,
а не независимым физическим эталоном.

Расхождения группируются по токовому диапазону,
несимметрии фаз, качеству измерений и источнику записи.
Отдельно рассматривается число переключений выхода.
Частое изменение решения может указывать на работу
около границы, на помеху либо на действительную
смену режима. Различить эти причины позволяет
возврат к сигналам и обстоятельствам события.

Для описания состава массива применяются также
многомерные признаки осциллограммы.
Они включают уровни токов и напряжений,
несимметрию, пик-факторы и временные характеристики.
Проекции главных компонент показывают общие различия,
группировка выделяет похожие записи, а поиск
изолированных примеров помогает находить редкие режимы.
Такие методы направляют предметную проверку.
Положение точки на проекции само по себе
не определяет правильное состояние защиты.

Результатом этого этапа становится не только сводная
таблица согласия, но и перечень физических областей,
в которых изменение органа наиболее существенно.
Это одинаково полезно при пересмотре аналитической
характеристики, адаптивного порога и нейросетевого решения.

\iffalse
(Примечание автора: Много текста тут, а картинки или таблицы все дальше? И в статье? Прост как-то сухой припиской / началом другого это кажется. Может оно и так, но повторюсь, потерял в начале смысл, зачем каждое из этого?.)
(Примечание автора: Ок, ждём последующего сжатия текста, пока не приступаем к этому, думаю ещё рецензий на сам текст получить.)
\fi
\authorreviewnote{AA159}

Пик-фактор сопоставляет наибольшую мгновенную величину сигнала
с его действующим значением. Он помогает выделить записи
с кратковременными выбросами. Несимметрия и временные признаки
дополняют этот отбор, чтобы при просмотре были представлены
разные по характеру процессы, а не только разные уровни тока.

\iffalse
(Примечание автора: Оно вроде бы много говорится, всякие пик факторы, а что это? А зачем это? В общем надо какую-то и общую понятность, и всё пояснять от вопроса «зачем мы это делаем?».)
(Примечание автора: Ок, закрываем.)
\fi
\authorreviewnote{AA160}
